\sectionguard
\section*{Source-Grounded Synthesis Across the Propers}

Nothing in this section is new evidence. It states what the checked witnesses
and the collated texts together support once the ten elements are read as one
book-order, and it stops where they stop. Proposals that go further are in
\textit{Interpretive Possibilities}, and the boundary is not decorative: three
of the four relations below were reached by comparing sources that disagree.

\subsection*{1. The same Latin word fails twice, and the second failure is a story}

The Epistle's technical term for what Moses' covenant was is
\latin{ministrátio}: \latin{ministrátio mortis}, the ministration of death
engraved in letters on stones, set against \latin{ministrátio spíritus}. The
Gospel then puts two ministers on a road. Augustine's table on the parable,
verified in the Latin, reads priest and Levite as the \latin{sacerdotium et
ministerium Veteris Testamenti, quod non poterat prodesse ad salutem} --- the
priesthood and ministry of the Old Testament, which could not profit unto
salvation. That is not a resemblance noticed here; it is the same noun and the
same judgment, made by a witness expounding the Gospel and not thinking about
the Epistle. Paul argues the point in the abstract, with tables and glory and
an a fortiori; Luke stages it as two men who look and walk on. The formulary
prints them forty lines apart.

What that convergence does \textit{not} license is worth saying, because the
tradition itself refused it. Ambrosiaster reads \latin{líttera enim occídit}
juridically --- the written law was given to Moses \latin{ut contemnentes Legem
occideret}, to execute its despisers --- and Chrysostom the same way: the Law
seizes a murderer and puts him to death, the Gospel seizes a murderer and
gives him life. On that reading the letter kills the guilty, and the Old
Testament's ministry is deficient rather than culpable. The older Augustine
moves the killing inward: the letter is any command without grace, the
Decalogue included, because the bare precept shows what to do and supplies no
power to do it. Aquinas then formalises Augustine's own definition of
``letter'' and draws the conclusion that closes every exit ---
\latin{littera etiam Evangelii occideret, nisi adesset interius gratia fidei
sanans}. Held together, these witnesses forbid the cheapest reading of this
Sunday: that the priest and the Levite failed because they were Jewish
ministers, and that the Christian reader is on the Samaritan's side by
membership. On Augustine's and Aquinas's account the letter that cannot heal
is the reader's own text, and the ministry that walks past is available to
anyone holding one.

\subsection*{2. Sufficiency is confessed before it is argued}

The Collect concedes the Epistle's thesis in its opening subordinate clause,
before it asks for anything: \latin{de cuius múnere venit, ut tibi a fidélibus
tuis digne et laudabíliter serviátur}. Worthy service is God's gift, not the
petitioner's contribution. The Epistle then supplies the argument the oration
took as given --- \latin{non quod sufficiéntes simus cogitáre áliquid a nobis,
quasi ex nobis \dots\ sed sufficiéntia nostra ex Deo est} --- disclaiming not
merely the deed but the interior act. The link is not this guide's: Schuster's
commentary on this Sunday quotes that verse of 2 Corinthians against any
suspicion that the Collect overstates. And the Introit has already put the
same confession in the mouth of a body that cannot help itself, with the
tradition's most concentrated commentary on why: Cassian hands over this one
verse because ``it contains the thought of one's own weakness, confidence in
the answer, and the assurance of a present and ever ready help''; Augustine
grounds it at \latin{ego vero egenus et pauper sum}, ``help I seek, always
mine infirmity I confess.''

Three elements therefore make one claim in three grammars --- sung petition,
oration concession, apostolic argument --- and the transmission records show
the concession is not a late refinement: the Collect stands in the Verona
collection and in the Gelasian, with the 1962 wording exact in the Hadrianum,
and the postconciliar Missal still says it. What cannot be claimed is that
anyone assembled the three for this reason. The Hadrianum witness shows the
chants and the orations already travelling together for this Sunday; it does
not show why.

\subsection*{3. The chants answer in the voice of the person the Gospel vindicates}

Between the two lessons the Mass sings the posture of someone whose standing
is received. The Gradual's \latin{áudiant mansuéti} drew from Cassiodorus a
pointed exclusion --- \latin{non dixit, lege docti, non ieiunántes, non
psalléntes, sed mansuéti} --- not the law-learned, not the fasters, not the
psalm-singers, but the meek are bidden rejoice. He was expounding a psalm, not
a lectionary; the Gospel that follows in this book stages a
\latin{legisperítus} who is exactly \latin{lege doctus}, answers his citation
exam perfectly, and is corrected by a Samaritan. Basil, on the same clause,
had already made the meek those who win ``without battle,'' and guessed
prophetically at the disciples --- whose eyes the same Gospel calls blessed.
The convergence is documented at both ends and belongs to the reception of two
separate texts; the liturgical juxtaposition is the missal's, and no compiler's
motive is asserted from it.

The Alleluia keeps the day honest about how long that posture lasts. Its verse
is the opening of the psalter's one lament that never turns, and the checked
witnesses divide on the speaker: Augustine, Cassiodorus and Bellarmine hear
Christ in the Passion, with the day-and-night pair concretised at last into
Good Friday and Gethsemane; Theodoret hears exilic Israel and, expressly,
``the pains of all human nature after sin.'' The two readings are not
reconciled here. Both leave the chant doing the same liturgical work: a cry
that does not stop, framed in \latin{Allelúia} --- unstopped petition sung
inside praise, which is precisely the difference between the Introit's plea
and despair.

\subsection*{4. Intercession is carried to the altar and answered in kind}

The Offertory is the formulary's hinge, and its evidence is unusually hard.
The chant is an Old-Latin, Septuagint-shaped adaptation of Exodus 32, and the
proof is checkable rather than stylistic: Augustine's own lemma in the
\work{Quaestiones in Heptateuchum} reads \latin{de malitia quam dixit facere
populo suo}, the chant's construction and not Jerome's. What the Church
therefore sings while the gifts are carried up is a text older than the Bible
printed beside it, and what it narrates is a man interceding for people who
have just disqualified themselves.

The reception of that scene is where the day's doctrine becomes an action
rather than a thesis. Augustine, on God's \latin{Sine me}, hears the whole
mechanism: the words are spoken \latin{ac si diceretur \dots\ dilectio tua in
illos intercedit mihi} --- your love intercedes with me for them --- so that
``when our own merits weigh us down, we can be relieved before him by the
merits of those he loves.'' Gregory reaches the same place from the other
side: to say ``let me alone'' to a servant is to give him boldness to plead,
and God ``bears the supplicant's contradiction which He inspires.'' Tertullian
alone among the checked witnesses names the figure outright --- Moses
deprecating the Father's wrath is ``Christ, prefigured \dots\ as the
deprecator of the Father, and the offerer of His own life for the salvation
of the people'' --- and that is worth marking as one witness rather than a
consensus, because preaching commonly attributes it to the tradition at large.

The two orations that follow convert intercession into exchange, and their
grammar carries the argument. The Secret's gerund makes God's act of lavishing
pardon the very mode in which the offerings honour his name --- \latin{ut,
nobis indulgéntiam largiéndo, tuo nómini dent honórem} --- so that
propitiation and praise, which a treatise would distinguish, are one clause.
The Postcommunion then opens on the Epistle's own verb: \latin{spíritus autem
vivíficat} in the lesson, \latin{Vivíficet nos} over the sacrament. That is a
textual observation about shared vocabulary on two collated texts, and nothing
more is claimed from it; but what the prayer asks with that verb answers the
Gospel's diagnosis exactly. \latin{Páriter} refuses to rank expiation and
\latin{munímen}: cleansing for what is past, a rampart for what is coming.
Augustine's \work{Sermo} 131 had said the same of the inn --- baptism forgives
all sins, and yet \latin{adhuc curatur}, the man is still under treatment ---
which is why the Church's last word on this Sunday is not thanksgiving for a
completed rescue but a request that the rescue keep working.

\subsection*{What the whole formulary supports, and where it stops}

Read in book order, the ten texts hold one position that no single element
states: that diagnosis and cure are different acts, and that this Mass is
arranged around the second. The letter can show the wound and cannot bind it;
the Law's own ministers see and pass; sufficiency, worthy service, rescue,
pardon and life are each asked for or confessed as received. Mercy is what
crosses the gap --- \latin{Non enim cognatio facit proximum, sed misericordia}
in Ambrose's sentence, and \latin{Vade, et tu fac simíliter} in the Lord's.

Two limits control that statement. First, it is a reading of the assembled
formulary and of checked reception, not a recovered intention: the one piece
of compositional evidence in hand, the Hadrianum's cue line pairing
\latin{Deus in adiutorium}, \latin{Benedicam dominum}, \latin{Precatus est
moyses} and \latin{De fructu operum}, proves early association and says
nothing about why. Second, Bede's rule governs the allegory throughout --- the
parable specially designates the Son of God \textit{and} teaches that everyone
who shows mercy is a neighbour, and the rules of allegory must not
\latin{moralia mutuae fraternitatis instituta \dots\ extenuare et auferre}.
A synthesis that made this Sunday only a doctrine of grace would fall under
Bede's own prohibition. The imperative is the last thing the Gospel says.
